% ======================================================================
%  TIER 2 KARTLARI: Understanding Users
% ======================================================================

\card{Intent}{kullanıcının gerçekte ne istediği}
\cl{Problem} Kullanıcının \emph{yazdığı}, \emph{istediği} şey değildir.
``Fatura'' diyen biri onu görüntülemek mi, itiraz etmek mi, ödemek mi istiyor?
\cl{Tanım} Intent, bir ifadenin ardındaki gerçek amaç --- yüzeydeki kelimeler değil.
\cl{Model} İyi bir garson siparişi değil, müşterinin \emph{açlığını} anlamaya çalışır.
\cl{Ne zaman} Çok amaçlı asistanlarda; doğru akışa veya araca yönlendirmeden önce.
\cl{Tuzak} Yüzeydeki kelimelere göre dallanmak; niyeti çıkarmadan yanıtlamak.
\cl{İlişkili} \term{Prompt} · \term{Router} \,{\scriptsize\color{softgray}(İlke \#22)}

\card{Conversation \& Persona}{çok turlu durum ve modelin kimliği}
\cl{Problem} Tek mesaj değil, turlar arası bir durum vardır --- ve modelin tutarlı
bir kimliği olmalıdır. Aksi hâlde her turda farklı bir ``kişi'' cevap verir.
\cl{Tanım} \term{Conversation} turlar arası durumdur; \term{Persona} modelin
tutarlı rol, ton ve kimliğidir (system prompt'ta tanımlanır).
\cl{Model} İyi bir telefon operatörü: her çağrıda aynı kişilik, ama bağlamı hatırlar.
\cl{Ne zaman} Çok turlu ürünlerde; marka tonu ve tutarlılık önemliyse.
\cl{Tuzak} Her turda farklı kişilik; ya da tüm geçmişi context'e yığıp tutarlılığı bozmak.
\cl{İlişkili} \term{System Prompt} · \term{Memory} · \term{Context Engineering} \,{\scriptsize\color{softgray}(İlke \#2)}
